% Generated draft: requires mathematical and source review.
\documentclass[11pt]{article}
\usepackage[margin=25mm]{geometry}
\usepackage{fontspec}
\setmainfont{Latin Modern Roman}
\usepackage{amsmath,amssymb,mathtools,unicode-math}
\setmathfont{Latin Modern Math}
\usepackage{enumitem,needspace,xurl,hyperref}
\setlength{\parindent}{0pt}
\setlength{\parskip}{5pt}
\setlength{\emergencystretch}{3em}
\newcommand{\N}{\mathbb{N}}
\newcommand{\Z}{\mathbb{Z}}
\newcommand{\Q}{\mathbb{Q}}
\newcommand{\R}{\mathbb{R}}
\newcommand{\C}{\mathbb{C}}
\newcommand{\F}{\mathbb{F}}
\newcommand{\A}{\mathbb{A}}
\newcommand{\PP}{\mathbb{P}}
\newcommand{\E}{\mathbb{E}}
\newcommand{\eps}{\varepsilon}
\newcommand{\dd}{\,\mathrm{d}}
\begin{document}
\section*{UnsolvedMath ID: 30001731}
\textbf{Source problem code:} catalogue entry
\subsection{Definitions and mathematical statement}
The names \emph{Plimpton 322} and \emph{MS 3972 no.\ 3} designate the historical texts named in the supplied record. The record describes the latter as analogous to the former, but supplies neither their contents nor a precise specification of the analogy. Here, ``tabular format'' refers to the presentation of information in Plimpton 322; no mathematical definition of that format is supplied.

The source asks: What is the mathematical meaning of the tabular format of Plimpton 322, and why is its information presented tabularly while the analogous text MS 3972 no.\ 3 is not?

This is a historical and interpretive research question. The supplied evidence does not specify mathematical objects, hypotheses, quantifiers, or criteria for a satisfactory interpretation sufficient to convert it into a precise mathematical proposition. Both the mathematical significance of the arrangement and the historical explanation of the difference in presentation remain the requested targets; neither is asserted here.
\subsection{Short English statement}
Explain the mathematical significance of the tabular presentation of Plimpton 322 and its contrast with the text identified as MS 3972 no.\ 3. The supplied record frames this as an interpretive historical question and does not provide enough evidence for a precise theorem-style formulation.
\subsection{Sources}
\begin{itemize}
\item Supplied local UnsolvedMath registry record, permanent ID \(30001731\), title \emph{Tabular Meaning of Plimpton 322}, source problem code \texttt{catalogue entry}, in the Oberwolfach Reports Open Problems collection. No individual proposer is identified; the supplied proposed year is 2011.
\item \emph{Mini-Workshop: History of Numerical and Graphical Tables} (2011), Oberwolfach Reports, pp.\ 639--689, DOI 10.4171/owr/2011/12. Supplied source report and literature reference labeled ``Oberwolfach Report 12/2011'': \url{https://doi.org/10.4171/owr/2011/12}. The supplied background describes Oberwolfach Reports as documenting mathematical workshops and their open-problem discussions.
\item The supplied literature assessment, checked on \texttt{2026-08-21}, labels the question ``open'' and characterizes the meaning of the tabular format in these historical examples as an interpretive research question, not a resolved mathematical proposition. This is supplied dated curation derived from the OWR record, not an independently verified literature assessment; time-sensitive status and cited claims require independent verification.
\item The registry supplies the statement-status label \texttt{corrected\_verified} and states that the official report confirms the historiographical question and that an extraction label was removed. These are supplied verification claims, not independent verification performed for this draft.
\end{itemize}
\end{document}
